\BeleskeID{04-s013}
% element: 04-s013:e01 — Tabelarni postupak 13:2 sa svih pet redova i naglašenim poslednjim nenultim korakom
% element: 04-s013:e02 — Tabelarni postupak 0.2×2 sa svih devet koraka i oznakom nastavka
% element: 04-s013:e03 — Obe strelice čitanja cifara i rezultat 13.2=1101.001100110…
% element: 04-s013:e04 — Tačniji konačan zapis, dogovor broja cifara i tačni primeri 0.5, 0.25, 0.75
% element: 04-s013:d01

Razlomljeni ostatak se posle četiri množenja vraća na početnu vrednost:
\[0.2\xrightarrow{\times2}0.4\xrightarrow{\times2}0.8
\xrightarrow{\times2}1.6\quad\text{(ostatak }0.6\text{)},\qquad
0.6\xrightarrow{\times2}1.2\quad\text{(ostatak }0.2\text{)}.\]
Izdvojene cifre su $0,0,1,1$. Isti ostatak znači da se ponavljaju isti koraci, pa je period razlomljenog zapisa $0011$. U koraku sa proizvodom $1.6$ nastavljamo sa $0.6$, a sa $1.2$ nastavljamo sa $0.2$.

Nijedan konačan odsečeni početak ovog periodičnog zapisa nije potpuno tačan. Potrebno je dogovoriti broj mesta; više mesta smanjuje grešku. Tačno konačne konverzije iz drugih primera su $0.5_{10}=0.1_2$, $0.25_{10}=0.01_2$ i $0.75_{10}=0.11_2$. Kod celobrojnog dela red $0{:}2$ proizvodi dodatnu vodeću nulu, pa nije nova značajna cifra.
